\section{Finite arithmetic and cyclic groups}

\subsection{Congruence and finite fields}

For an integer $n>1$, write $a\equiv b\pmod n$ when $n$ divides $a-b$.
The residue classes form the ring $\Z/n\Z$. Addition and multiplication are
performed normally and then reduced to a representative modulo $n$.

\begin{definition}[Finite prime field]
For a prime $p$, the field
\[
  \F_p=\Z/p\Z
\]
has $p$ elements. Every nonzero $a\in\F_p$ has a unique inverse $a^{-1}$
satisfying $aa^{-1}=1$.
\end{definition}

Primality matters. If $n=rs$ is composite, then the nonzero classes of $r$
and $s$ multiply to zero, so neither can have an inverse. Conversely, Bezout's
identity gives an inverse for every nonzero class modulo a prime. Thus
$\Z/n\Z$ is a field exactly when $n$ is prime.

Division in $\F_p$ means multiplication by an inverse. The extended Euclidean
algorithm computes that inverse, while Fermat's little theorem provides the
identity
\[
  a^{-1}=a^{p-2}\pmod p\qquad(a\ne0).
\]
The latter is particularly useful when an implementation already has an
efficient fixed-pattern exponentiation routine.

The nonzero elements form the multiplicative group
\[
  \F_p^\times=\F_p\setminus\{0\},\qquad |\F_p^\times|=p-1.
\]
Do not confuse this group with the field's additive group, whose identity is
$0$ rather than $1$.

\subsection{Groups, generators, and order}

\begin{definition}[Group]
A group $(\G,\circ)$ has an associative binary operation, an identity $e$,
and an inverse for every element. It is abelian when $x\circ y=y\circ x$.
\end{definition}

For $g\in\G$, the cyclic subgroup generated by $g$ is
\[
  \langle g\rangle=\{g^k:k\in\Z\}.
\]
Its size is the order $\ord(g)$, equivalently the smallest positive $r$ for
which $g^r=e$. In additive notation the same statements become
\[
  \langle G\rangle=\{[k]G:k\in\Z\},
  \qquad
  \ord(G)=\min\{r>0:[r]G=\OO\}.
\]
Square brackets show that scalar multiplication is repeated group addition;
it is not multiplication of the coordinates.

Lagrange's theorem says the order of a subgroup divides the order of the
finite ambient group. Consequently, cryptographic parameter design is driven
by the factorization of the group order, not merely by the number of bits in
the field modulus.

\paragraph{A running multiplicative example.}
In $\F_{29}^\times$, the element $2$ has order $28$ and generates the whole
group. Forward evaluation quickly gives
\[
  2^{19}\equiv26\pmod{29}.
\]
The inverse question ``which exponent produced $26$?'' is a tiny instance of
the discrete logarithm problem. The first marimo notebook makes the orbit and
all inverses visible.

\subsection{Quadratic residues}

A nonzero $r\in\F_p$ is a quadratic residue when $r=y^2$ for some $y$.
Euler's criterion packages this test as the Legendre symbol
\[
\left(\frac{r}{p}\right)=r^{(p-1)/2}\pmod p\in\{1,-1\}.
\]
Exactly half of the nonzero field elements are squares. Later, this one-bit
classification tells us whether a Montgomery $u$-coordinate lifts to a point
on a curve or on its quadratic twist.

\begin{takeaway}
Keep three integers separate: the field prime $p$, the order $\ell$ of the
chosen prime-order subgroup, and the cofactor $h$ satisfying
$\#\G=h\ell$.
\end{takeaway}
